% ================================================================
% chapters/parts/ch07-b.tex —— 第 7 章 原位红外光谱方法（分片 b）
% 覆盖：原书 §7.3（探针分子、金属催化剂表征）与 §7.4.2—7.4.4
%        （表面羟基、固体表面酸性、氧化物表面氧物种）
%        书内 p.448—488 / PDF p.467—507
% 注：原书 §7.2「吸附分子的特征和它的红外光谱诊释」属分片 a 的范围，
%     本节只从「判断」角度引用其结论，不重复叙述。
% 版式自查（WRITER_BRIEF §7.3）：单栏版心 165 mm；\onecolwidth = 0.80\textwidth ≈ 132 mm。
%   图 \ref{fig:ch07b-co-model} 用 groupplot，3 列：
%     3×48 + 2×8 = 160 mm，加左侧轴标签 ≈3 mm → ≈163 mm ≤ 165 mm  ✓
% ================================================================

\subsection{探针分子的选择与红外光谱诊释}

\begin{probout}
拿到一张吸附分子的红外谱，问题从来不是「有没有峰」，而是三个判断：峰在哪
（位置）、峰多高（强度）、峰随条件怎么动（位移方向与幅度）。三者分别对应
表面有什么位、位有多少、位上的电子状态如何。本节要输出的不是某个体系的
具体归属，而是一套可复用的判读路径：先选对探针分子，再用位置定类型、用
强度定数量、用位移定电子状态，最后回查各环节的干扰。
\end{probout}

\subsubsection{从谱带到表面位的判读链条}

\begin{principle}
吸附分子的红外光谱与它吸附前的光谱相比，除了可能出现新的「表面键」伸缩
振动谱带，原有振动频率也会位移；若吸附后分子化学结构改变（如双键打开），
振动改变更大\src{S1 §7.2，书内 p.444 / PDF p.463}。判读因此有两条依据：
\begin{itemize}
  \item \textbf{位置}定「位」的类型。同一探针分子在不同性质的表面位上以不同
        方式键合，其基团特征频率落在不同区间，故位置先给类型。
  \item \textbf{强度}定「位」的数量。定量只能建立在朗伯—比尔定律之上：
\end{itemize}
\begin{equation}
  A=\varepsilon\,c\,l
  \label{eq:ch07b-beer}
\end{equation}
式中 $A$ 为积分吸光度（谱带面积），$\varepsilon$ 为该谱带的积分消光系数，
$c$ 为吸光物种的表面浓度（单位面积上的摩尔数或分子数），$l$ 为有效光程。
若把探针分子近似为「一位一分子」的化学计量滴定，则 $c$ 即该类型表面位的
表面密度。

\textbf{成立条件（三条必须同时成立）}：一是化学计量数已知且不随组成变化；
二是各吸附构型（线式／桥式／多点配位）的消光系数已知且恒定；三是不同样品
的有效光程 $l$ 可比。三条中任意一条不成立，强度就只是相对强弱，不能再读
成位密度。
\end{principle}

\begin{measure}
探针分子的谱带位置本身对条件很敏感，测量时必须同时锁定：
(1) \textbf{温度}——吸附温度与抽空温度共同决定留下哪些物种，例如吡啶需在
$150\,\degC$ 抽空以排除物理吸附（见 \S\,\ref{sec:ch07b-oxide}）；
(2) \textbf{覆盖度}——同一吸附位的 $v(\mathrm{C{-}O})$ 会随覆盖度漂移，故对比样品
须在同一平衡压力与同一抽空时间下测；
(3) \textbf{样品形态}——自支撑片的厚度与散射本底直接进入方程
\equref{eq:ch07b-beer} 的 $l$，厚度不同的片强度不可直接比；
(4) \textbf{本底}——须对同一样品在同一条件下采集吸附前本底谱并作差谱。
\end{measure}

\subsubsection{常用探针分子及其适用范围}

选择探针分子的实质是选择它的\textbf{碱性（或酸性）强度}与\textbf{分子尺寸}
两个参数。碱性太强的探针会把弱酸位一并质子化，从而抹平强度差异；尺寸太大
的探针进不了小孔，只能报告外表面。表 \ref{tab:ch07b-probes} 给出本节与
\S\,\ref{sec:ch07b-oxide} 反复用到的五种探针。

\begin{table}[htbp]
\centering
\caption{常用探针分子的适用范围、谱带位置与干扰}
\label{tab:ch07b-probes}
\small
\begin{tabularx}{\onecolwidth}{@{}>{\raggedright\arraybackslash}p{13mm}
 >{\raggedright\arraybackslash}p{30mm}
 >{\raggedright\arraybackslash}p{29mm}
 >{\raggedright\arraybackslash}X
 >{\raggedright\arraybackslash}X@{}}
\toprule
探针 & 测什么 & 谱带位置区间$^{\mathrm{a}}$ & 定量前提 & 主要干扰 \\
\midrule
CO &
 金属位（顶位／桥位）、氧化物表面 L 酸位与配位不饱和阳离子；
 \textbf{不与 B 酸位作用} &
 顶位／L 酸位配位 $>2000$；桥位 $<2000$；
 氧化物上 L 酸位约 $2150$—$2200$ &
 需各类吸附构型的积分消光系数；
 需按覆盖度归一 &
 气相 CO 的振动—转动谱带（$P$／$R$ 支）；
 桥位与顶位谱带重叠；
 CO 在金属上可歧化生成 $\ce{CO2}$ 与表面碳 \\
吡啶 &
 L 酸位与 B 酸位的同时区分（环变形振动对成键方式敏感） &
 B 酸位 $1540$、L 酸位 $1450$（均为环变形振动）；
 另有 $1490$ 带为两类位共有 &
 需 $150\,\degC$ 抽空扣除物理吸附与气相；
 需分别标定两类位的消光系数 &
 物理吸附吡啶（室温可抽除）残留；
 探针在强酸位上的质子化产物本身即被测物种 \\
$\ce{NH3}$ &
 L 酸位与 B 酸位（主要以变形振动区分） &
 L 酸位配位 $\ce{NH3}$：$3330$、$1610$；
 B 酸位 $\ce{NH4+}$：$3230$、$1432$ &
 碱度大，滴定可近化学计量；
 但强弱酸位同时被滴定 &
 \textbf{碱性过强}，会把弱酸位一并滴定而掩盖其差异；
 易与表面羟基形成氢键物种，产生额外谱带 \\
$\ce{CO2}$ &
 表面碱性位（碱性氧位）、表面碳酸盐与甲酸盐物种 &
 碳酸盐／羧酸盐的 $v_{\mathrm{as}}(\mathrm{OCO})$ 约 $1572$、
 $v_{\mathrm{s}}(\mathrm{OCO})$ 约 $1369$；
 甲酸盐 $\nu(\mathrm{C{-}H})$ 约 $2870$、$\delta(\mathrm{C{-}H})$ 约 $1379$ &
 需先扣除表面预存碳酸盐保护层；
 需已知各碳酸盐构型的消光系数 &
 碱性氧化物表面天然覆盖一层稳定碳酸盐；
 多种碳酸盐构型谱带重叠 \\
$\ce{CH3CN}$ &
 表面羟基的酸性强度（由氢键引起的羟基位移反推） &
 以羟基伸缩带位移读数：二氧化硅—氧化铝上
 $3750$ 一带吸附后分裂为两个峰 &
 需记录吸附前后羟基带的半宽与位移；
 需同时测 OD 形式作对照 &
 探针在强酸位上可被质子化，生成物种的谱带与氢键物种不同；
 液态／物理吸附态的谱带覆盖 \\
\bottomrule
\end{tabularx}

\vspace{1pt}
{\srcsize $^{\mathrm{a}}$ 单位 $\mathrm{cm^{-1}}$。CO 在氧化物 L 酸位上的
$2150$—$2200$ 区间与在金属顶位上的 $>2000$ 区间在数值上有重叠，
两者必须靠样品体系（有无金属组元）与共吸附实验区分，不能只凭波数。
CO 位置的区间系按本书所给实例归纳，具体体系的数值须逐例核对。}
\end{table}

\begin{dataflow}
从一张吸附谱到一个结论，应当走完四步：
\textbf{(1) 定本底与差谱}：取得同一样品吸附前谱，作差得吸附物种谱。
\textbf{(2) 分峰定位置}：对重叠区作分峰，读出各组分的位置与面积；只有位置
确定后，才谈得上归属。
\textbf{(3) 用位置定类型}：按 \tabref{tab:ch07b-probes} 的区间先给「哪一类位」，
再在同类位内部用共吸附或同位素实验细分。
\textbf{(4) 用面积与位移定量}：面积经方程 \equref{eq:ch07b-beer} 换算成表面
浓度，位移用于判读电子转移方向与程度。
四步中任何一步的中间结论都不得越过下一步直接当结论用——尤其是「位置」
只能给类型，不能给强度。
\end{dataflow}

\begin{pitfall}
\textbf{误判一：把谱带位移作单因素归因。}同一吸附位的 $v(\mathrm{C{-}O})$
位移至少有四个来源：吸附构型改变（顶位／桥位／多点配位）、覆盖度变化、
金属组元的电子效应、以及偶极—偶极耦合造成的集体位移。本书的数据正好
把前两者的量级暴露出来：$\mathrm{Pd{-}Ag}$ 体系中覆盖度从高到低变化引起的
CO 化学位移只有约 $9\;\mathrm{cm^{-1}}$，而合金组成变化引起的位移达
$29\;\mathrm{cm^{-1}}$\src{S1 §7.3.2，书内 p.453 / PDF p.472}。因此
「位移大就是电子效应」在量级上可以成立，但\textbf{必须先给出几何效应
与覆盖度效应的上限估计}，否则无法排除它们。

\textbf{误判二：忽略消光系数差异而直接比强度。}顶位与桥式 CO 的积分消光
系数不同，$\ce{NH3}$ 在 L 酸位与 B 酸位上的变形振动消光系数也不同。
若拿两类谱带的面积直接相比，得到的比值中混入了消光系数的比值。

\textbf{误判三：忽略探针分子自身的反应。}CO 在金属表面上可发生歧化，
生成的 $\ce{CO2}$ 与表面碳会同时改变谱图与表面状态；吡啶在强酸位上会
质子化，其质子化产物恰恰是被测物种，二者在谱上不易分开。

\textbf{误判四：自支撑片厚度不一而直接比较强度。}不同样品压片厚度不同，
方程 \equref{eq:ch07b-beer} 中的 $l$ 就不同。凡跨样品比较强度，必须先归一：
常用做法是以分子筛骨架振动或二氧化硅的 $1860\;\mathrm{cm^{-1}}$ 附近
（二氧化硅骨架的倍频／组频）谱带为内标，或直接称量压片并记录厚度。
\end{pitfall}

\subsubsection{一个读图实例：从「峰在哪」到「位是什么」}

下面的实例取自本书对 $\ce{Pt{-}Ru/SiO2}$ 双金属催化剂的双探针研究
\src{S1 §7.3.1，书内 p.449—451 / PDF p.468—470}。请把它当作判读路径的
示范，而不是一个可以直接背下来的归属表。

\begin{obs}
在含 $38\%$ Ru 的 $\ce{Pt{-}Ru/SiO2}$ 上，CO 与 NO 混合气于室温竞争吸附，
抽空后测得：CO 吸附峰约在 $2065\;\mathrm{cm^{-1}}$，NO 吸附峰约在
$1550\;\mathrm{cm^{-1}}$，三个峰在室温抽空后都是稳定的。
随样品中 Ru 原子百分数增加（$10\%\to100\%$）：NO 峰相对 CO 峰的强度
增大；NO 吸收峰向高波数位移；CO 谱带强度减弱；CO 谱带向低波数位移，
对应位置由约 $2074\;\mathrm{cm^{-1}}$ 移到约 $2030\;\mathrm{cm^{-1}}$
\src{S1 表 7-（Pt-Ru 吸附数据），书内 p.450 / PDF p.469}。
\end{obs}

\begin{infer}
由「CO 峰在 $2065\;\mathrm{cm^{-1}}$ 以上、NO 峰在 $1550\;\mathrm{cm^{-1}}$
附近」这两个位置，结合两种探针在各金属上的成键差异，可推断 CO 主要吸附
在 Pt 中心上、NO 主要吸附在 Ru 中心上。这一步只用到了\textbf{位置}，
因此它只给出「哪种分子占哪类中心」，不给任何数量结论。

由「NO 峰强度随 Ru 含量增大而相对增强、CO 峰随之减弱」这一\textbf{强度}
变化，可进一步推断两种探针确实分别报告了两个组元的表面分数——这是把
强度当作表面组成的量度来用的前提。

由「CO 谱带随 Ru 含量增加而向低波数位移」这一\textbf{方向}，按
$d$—$\pi^{*}$ 反馈模型可推断：金属向 CO 的 $\pi^{*}$ 反键轨道反馈电子
的程度增加，$v(\mathrm{C{-}O})$ 因而降低。这一推断的强度弱于前两条：
位移方向支持电子转移的判断，但位移的\textbf{幅度}不能直接当作电子转移量，
理由见上文误判一。
\end{infer}

\begin{crosscheck}
三条互补的回查：
(1) 用已知分散度的单一金属样品（$\ce{Pt/Al2O3}$、$\ce{Ru/Al2O3}$）分别
标定两种探针在各自金属上的谱带位置与积分消光系数，再把同一套条件用到
双金属样品上；
(2) 用化学吸附（$\ce{H2}$、CO 吸附量）或 TPD 给出的表面总原子数与红外
定量结果对照，二者相符才说明「位置—强度」链条没有系统性偏差；
(3) 用同位素或共吸附实验独立验证电子转移方向，例如改变共吸附分子的
给电子能力，看位移方向是否随之翻转。
本书对 $\ce{Pt{-}Ru}$ 体系的结论正是靠化学吸附对照得到支持的
\src{S1 §7.3.1，书内 p.451 / PDF p.470}。
\end{crosscheck}

\subsection{红外光谱在金属催化剂表征中的应用}

\subsubsection{双探针法测定双金属催化剂的表面组成}

\begin{probout}
要输出的是\textbf{最表面层}的组成，而不是体相组成。合金催化剂的表面组成
常与体相明显不同，催化性能随之显著变化。二次离子质谱、离子散射谱、
俄歇电子能谱等表面方法设备昂贵，且所测数据不都是最表面层，不易与催化
性能直接关联；常规化学吸附法测定双金属表面组成又仅限于第 VIII 族与
第 IB 族组成的合金\src{S1 §7.3.1，书内 p.449 / PDF p.468}。双探针红外
方法要补的正是这个缺口：用两种选择性地吸附在不同组元上的探针分子，
以它们的谱带强度比读出表面组成。
\end{probout}

\begin{principle}
设两种探针分别只吸附在两个组元上，且各自的积分消光系数 $K$ 与吸附构形
系数 $S$ 都与表面组成无关，则两条谱带的积分吸光度分别正比于各自的表面
浓度：
\begin{align}
  A_{\mathrm{CO}} &= K_{\mathrm{CO}}\,S_{\mathrm{CO}}(\mathrm{Pt})\,x_{s},
  \label{eq:ch07b-co-abs}\\
  A_{\mathrm{NO}} &= K_{\mathrm{NO}}\,S_{\mathrm{NO}}(\mathrm{Ru})\,(1-x_{s}),
  \label{eq:ch07b-no-abs}
\end{align}
式中 $x_{s}$ 为 Pt 在\textbf{表面}的原子分数，$1-x_{s}$ 为 Ru 的表面原子
分数，二者按归一化定义。
两式相除消去共同因子，得
\begin{equation}
  \frac{A_{\mathrm{CO}}}{A_{\mathrm{NO}}}
  =\frac{K_{\mathrm{CO}}S_{\mathrm{CO}}(\mathrm{Pt})}{K_{\mathrm{NO}}S_{\mathrm{NO}}(\mathrm{Ru})}
  \cdot\frac{x_{s}}{1-x_{s}} .
  \label{eq:ch07b-ratio}
\end{equation}
若两种探针的消光系数与构形系数之比在整段组成范围内近似为常数，则
$A_{\mathrm{CO}}/A_{\mathrm{NO}}$ 对表面组成比 $x_{s}/(1-x_{s})$ 作图应为
直线。本书把 $A_{\mathrm{CO}}/A_{\mathrm{NO}}$ 对体相组成原子百分数
$(\mathrm{Pt})_{b}$、$(\mathrm{Ru})_{b}$ 作图，在相当宽的区间内得到直线
关系，据此判断在所研究体系内表面组成与体相组成一致
\src{S1 §7.3.1，书内 p.451 / PDF p.470}。
\tocheck{未对照原页面图像：式~\eqref{eq:ch07b-co-abs}、\eqref{eq:ch07b-no-abs}
在源 OCR 中作「$A_{\mathrm{CO}}=K_{\mathrm{CO}}\cdot S_{\mathrm{CO}}(\mathrm{Pt})_{s}$」
一类残形，表面分数下标的写法与是否含归一化因子无法从文本层判定；本节按
「表面分数」重写，符号约定须核原页 p.451。}
\end{principle}

\begin{measure}
该方法的实验条件须成套固定，缺一项比值就不可比：先在约 $\qty{6.67}{kPa}$
（$50\;\mathrm{Torr}$）CO 气氛下吸附至平衡并抽空，再与过量 NO 吸附平衡
后于室温抽空，最后重新暴露于 CO 气氛建立稳定覆盖
\src{S1 §7.3.1，书内 p.449—450 / PDF p.468—469}。
\tocheck{未对照原页面图像：上述三步的先后顺序与每步抽空时间在源 OCR 中
被截断（残存「空巧min」「抽空巧mln」等形），本节按「CO 平衡—抽空—NO
平衡—抽空—CO 再暴露」重建，抽空时长须核原页 p.449。}
记录的是室温抽空后仍稳定的谱带，因此只计化学吸附物种。
\end{measure}

\begin{dataflow}
顺序不可颠倒：\textbf{(1)} 在单金属参比样品上标定 $K$ 与 $S$；
\textbf{(2)} 在同一组条件下测一系列双金属样品的两条谱带面积；
\textbf{(3)} 计算比值并作图，检查线性区间；\textbf{(4)} 只有在线性成立的
区间内，才把比值换算成表面组成；\textbf{(5)} 用化学吸附给出的表面总原子数
交叉验证绝对值。本书明确指出该方法的两点优势——测得的是最表面层组成、
表面没有发生畸变——同时也留下三个待改进之处：定量精度还不够高、理论
消光系数变化规律不十分清楚、化学计量数随组成的变化规律有待研究
\src{S1 §7.3.1，书内 p.451 / PDF p.470}。
\end{dataflow}

\subsubsection{几何效应与电子效应的区分}

\begin{principle}
第二金属组元引入后改变吸附与催化性能，通常归因于两种效应：
\textbf{几何效应}（活性组元的原子被稀释，成对或多中心吸附位减少）与
\textbf{电子效应}（组元间发生电荷转移，改变金属的 $d$ 电子密度，
从而改变反馈键强度）。$d$—$\pi^{*}$ 反馈模型给出的判据是方向性的：
金属 $d$ 电子密度增加 $\Rightarrow$ 向 CO 的 $\pi^{*}$ 反馈增加
$\Rightarrow$ $v(\mathrm{C{-}O})$ 红移（向低波数）。
连接微观量与观测量的是方程 \equref{eq:ch07b-beer} 形式的
「强度—浓度」关系，而位置—电子密度之间\textbf{没有}可直接标定的
定量关系式，这正是下面必须靠量级比较来区分两种效应的原因。
\end{principle}

\begin{obs}
\textbf{$\mathrm{Pd{-}Ag}$ 体系}（总金属含量 $9\%$，合金物相与组成由 X 射线
衍射测定）：CO 在 Pd 上，高于 $2000\;\mathrm{cm^{-1}}$ 的谱带属顶位（线式）
吸附，低于 $2000\;\mathrm{cm^{-1}}$ 的属桥式吸附。Ag 含量增加时，桥式 CO
谱带强度明显下降直至完全消失，线式 CO 谱带强度明显增加
\src{S1 §7.3.2，书内 p.452 / PDF p.471}。
后续用傅里叶红外仪重新研究该体系，除强度规律一致外，还发现 CO 吸收带
随合金组成变化发生位移，Ag 含量增加时红移加大
\src{S1 §7.3.2，书内 p.452—453 / PDF p.471—472}。

\textbf{$\mathrm{Ru{-}Cu/SiO2}$ 体系}：Ru 上 CO 的谱带随 $\mathrm{Cu/Ru}$
原子比增加而红移，同时高频带的相对丰度减少；室温吸附 $\ce{H2}$ 后作 TPD，
$\ce{H2}$ 脱附温度随 $\mathrm{Cu/Ru}$ 原子比增加而下降
\src{S1 §7.3.2，书内 p.455 / PDF p.474}。
\end{obs}

\begin{infer}
\textbf{几何效应的判据}：Ag 对 Pd 起稀释作用，成双存在的 Pd 浓度随 Ag
含量增加而减少，因而桥式 CO 减少、线式 CO 增加。这一条只用到「强度」，
故它单独就能支持几何效应。

\textbf{几何效应不足的判据}：后来在 $\mathrm{Cu{-}Ni}$ 体系中进一步实验
发现两点——线式与桥式 CO 吸附态的比值变化，与用系统方法算出的双金属
原子对（如 Ni、Pd）表面浓度随合金组成的变化\textbf{不一致}；且由合金化
引起的谱带化学位移\textbf{大于} CO 覆盖度变化引起的化学位移。两点都指向
几何效应之外还有一个来源。

\textbf{电子效应的判据}：Ag 含量增加使 CO 吸附带红移加大，按反馈模型
说明由金属反馈到 CO 分子的电子增加，即 Pd—Ag 之间存在电子效应。本书
据此结论：$\mathrm{Pd{-}Ag/SiO2}$ 体系中几何效应与电子效应同时存在，
$\mathrm{Ni{-}Cu}$ 体系有类似现象\src{S1 §7.3.2，书内 p.455 / PDF p.474}。

\textbf{必须写明的一点}：这一推断的落脚点是「位移幅度大到几何效应解释
不了」，而不是「位移方向即电子效应的量度」。判据是\textbf{幅度比较}，
所以必须先取得覆盖度效应的量级上限。
\end{infer}

\begin{pitfall}
\textbf{把「位移大」直接读成「电子效应强」。}本书给出了两个量级：
覆盖度变化引起的 CO 化学位移只有约 $9\;\mathrm{cm^{-1}}$
（$2062\to2053$），而组成变化引起的位移为 $29\;\mathrm{cm^{-1}}$
（$2075\to2049$）\src{S1 §7.3.2，书内 p.453 / PDF p.472}。
两数相差三倍以上，才使「几何效应不足以解释」这一论断站得住。反之，
若某体系的组成位移只有十几波数，同一论断就不成立。

\textbf{把强度比直接当表面浓度比。}几何效应本身就改变构型分布，而不同
构型的消光系数不同，故强度比中同时含有「位有多少」与「构型怎么变」
两个信息。这也是本书把「理论消光系数变化规律不十分清楚」列为待改进项
的原因。

\textbf{忽略载体与预处理的影响。}$\mathrm{Ru{-}Cu}$ 体系的结果表明，
Cu 与 Ru 之间相互作用的强弱和电子转移的方向受载体性质及预处理条件影响；
体相 Ru 与 Cu 不相混溶，但分散在 $\ce{SiO2}$ 表面时可以明显相互作用
\src{S1 §7.3.2，书内 p.455 / PDF p.474}。因此同一对金属在不同载体上的
结论不可互推。

\textbf{忽略探针分子自身的反应与吸附位竞争。}CO 与 NO 共吸附时两者会
争夺同一类中心，若两种探针并非严格分区吸附，方程
\equref{eq:ch07b-ratio} 的分区前提就失效。
\end{pitfall}

\subsubsection{吸附分子间相互作用：把 $v(\mathrm{C{-}O})$ 当电子探针}

\begin{principle}
把这套判读用到「共吸附」上，CO 就不再只是被测对象，而成为报告表面电子
状态的探针。依据是反馈键在 CO 结合能中占很大份额——本书引用计算结果，
CO 在 Ni 上吸附的结合能有约 $84\%$ 来自反馈键贡献
\src{S1 §7.3.3，书内 p.456 / PDF p.475}。反馈键强弱与金属 $d$ 轨道情况
密切相关，因此\textbf{当共吸附分子的给电子能力不同时，$v(\mathrm{C{-}O})$
的位移方向与大小就报告了电子转移的方向与程度}。
\end{principle}

\begin{obs}
CO 预吸附在 Pt 上（$v_{\mathrm{CO}}=2065\;\mathrm{cm^{-1}}$），再引入
不同给电子能力的共吸附分子，测得的位移如下（电离势 $I$ 越小表示越容易
给出电子）\src{S1 §7.3.3，书内 p.456—457 / PDF p.475—476}：
\begin{itemize}
  \item $\ce{H2O}$（$I=12.6\;\mathrm{eV}$）共吸附：$v_{\mathrm{CO}}$ 由
        $2065$ 位移到 $2050\;\mathrm{cm^{-1}}$；
  \item $\ce{NH3}$（$I=10.5\;\mathrm{eV}$）共吸附：继续向低波数位移；
  \item 吡啶 $\ce{C5H5N}$（$I=9.2\;\mathrm{eV}$）共吸附：位移到
        $1990\;\mathrm{cm^{-1}}$；
  \item 受主型分子共吸附时方向相反：$\ce{HCl}$（$I=12.5\;\mathrm{eV}$）
        共吸附使 $v_{\mathrm{CO}}$ 位移到 $2075\;\mathrm{cm^{-1}}$，
        另一受主体系则出现 $2131\;\mathrm{cm^{-1}}$ 的谱带。
\end{itemize}
另有独立一例：$\ce{C6H6}$ 吸附在 $\ce{Pt/Al2O3}$ 上时，预吸附 CO 的
$v_{\mathrm{CO}}$ 由 $2065\;\mathrm{cm^{-1}}$ 红移
\src{S1 §7.3.3，书内 p.457 / PDF p.476}。
\tocheck{未对照原页面图像：上述各体系位移后的确切波数与
$\ce{NH3}$、$\ce{HCl}$、受主体系三组的具体数值在源 OCR 中残缺
（残存「位移到 cm」「$2131\;\mathrm{cm^{-1}}$（O—C—Pt—O）」等形），
本节只保留可判读的方向与能读出的数值，完整数值须核原页 p.456—457。}
\end{obs}

\begin{infer}
按反馈模型，给电子能力强的分子向 Pt 提供电子、使 Pt 的 $d$ 电子密度
增加，反馈程度加强，故 $v_{\mathrm{CO}}$ 向低波数位移；受主分子抽走电子，
反馈程度减弱，故向高波数位移。把五位共吸附分子按电离势排序，与实测位移
的排序一致——$\ce{H2O}<\ce{NH3}<$ 吡啶的给电子能力递增，红移也递增，
吡啶因最容易给出电子而得到最大的红移。\textbf{这是一条「位移报告电子
转移」的有效性论证}：结论落在「排序一致性」上，而非任何单一体系的
绝对位移值。

同一逻辑反过来用于判断吸附态的成键方式。苯在 Pt 上加氢与苯—氘交换的
机理长期存在解离吸附模型、非解离吸附模型与 $\pi$ 络合物模型之争；由于
这些模型都靠 $\mathrm{C{-}H}$ 键振动来判别，只能区分双键是否打开，
难以区分电子转移方向，故 $\pi$ 络合物模型长期未获实验证明。用 CO 作
探针后：吸附 CO 后再吸附 $\ce{C6H6}$，$v_{\mathrm{C{-}H}}<3000\;\mathrm{cm^{-1}}$
呈现饱和烃 $\mathrm{C{-}H}$ 特征，同时 $v_{\mathrm{CO}}$ 由
$2065\;\mathrm{cm^{-1}}$ 红移。按反馈原理，红移说明电子由 $\ce{C6H6}$
转移到 Pt 上，并在 Pt 表面形成 $\pi$ 络合物吸附态；加氢后这一作用明显
降低\src{S1 §7.3.3，书内 p.457 / PDF p.476}。
\end{infer}

\begin{crosscheck}
\textbf{用模型计算给判据定量。}位置—电子密度之间没有可直接标定的解析式，
故用模型计算把「$d$ 电子密度增加 $\to$ 红移」这条定性规律转成可比较的
曲线，再与实验点对照。\figref{fig:ch07b-co-model} 给出本节构造的模型
计算：以顶位吸附的 $v_{\mathrm{CO}}$ 为基准，按每个电子（以 $e$ 为单位）
约 $100\;\mathrm{cm^{-1}}$ 的红移率生成位移曲线，并按峰形叠加生成模拟谱。
图中同时标出本书实测的两类参照点——合金化引起的 $29\;\mathrm{cm^{-1}}$
组成位移与约 $9\;\mathrm{cm^{-1}}$ 的覆盖度位移——读者据此判断：在
$0.2$—$0.3$ 个电子的转移量级上，组成位移与覆盖度位移才可能被区分开。
\begin{figure}[htbp]
\centering
\begin{tikzpicture}
\tikzset{declare function={%
  % 参数名只用单字母，且避开 LaTeX 重音命令（\u、\x、\a 之外的单字母需自查）：
  % 用多字母名（如 \xx）会报 "Missing control sequence inserted"；
  % 用 \u、\x 会被当重音命令展开，报 "ill-formatted floating point number"。
  lgk(\p,\c,\a,\w)=(\a/(1+((\p-\c)/\w)*((\p-\c)/\w)));%
  mtop(\p)=2074-(100*\p);%
  mbr(\p)=1998-(100*\p);}}
\begin{groupplot}[
  group style={group size=3 by 1, horizontal sep=8mm,
               ylabels at=edge left, xlabels at=edge bottom},
  width=48mm, height=5.0cm,
  domain=1700:2200, samples=161,
  xmin=1700, xmax=2200,
  axis lines=left, xmajorgrids=false,
  tick align=outside,
  xlabel={波数$/\mathrm{cm^{-1}}$},
  xlabel style={font={\fontsize{8}{9.5}\selectfont}},
  x tick label style={/pgf/number format/1000 sep={},
                      font={\fontsize{7}{8.5}\selectfont}},
  y tick label style={font={\fontsize{7}{8.5}\selectfont}},
  ylabel style={font={\fontsize{8}{9.5}\selectfont}},
  legend style={font={\fontsize{7}{8.5}\selectfont}, draw=none, fill=none,
                at={(0.03,0.97)}, anchor=north west, cells={anchor=west}},
  ]
  \nextgroupplot[ymin=0, ymax=1.25, ytick=\empty,
                 ylabel={相对吸光度（模型）},
                 title={(a) 吸附构型}, title style={font={\fontsize{8}{9.5}\selectfont}}]
    \addplot[mainline, black]      {lgk(x,2074,1.00,16)};
    \addlegendentry{顶位 $2074$}
    \addplot[thinline, black, dashed] {lgk(x,1884,0.40,24)};
    \addlegendentry{桥位 $1884$（$\times2.5$）}
    \draw[gray] (axis cs:2000,0) -- (axis cs:2000,1.15);
    % 【重叠修正，2026-09-30】$2000$ 原用 anchor=south east 挂在 (1996,1.05)，
    % 与图例第一行 `顶位 2074` 完全压在一起：实测 `位`×`0` 墨迹相交 3.5×4.7 pt、
    % 轮廓交叠 1.078 pt²，另有 `位`×`2`、`2`×`0` 两对，合计 3 对。
    % 图例（实测 x 174—246、y 424—446）压不下来，改把标注移到竖线右侧、
    % 图例下方（该处两条谱线均接近 0，实测空白）。
    \node[anchor=west, font={\fontsize{7}{8.5}\selectfont}]
      at (axis cs:2008,0.60) {$2000$};
  \nextgroupplot[ymin=0, ymax=1.25, ytick=\empty,
                 title={(b) 电子密度（几何不变）},
                 title style={font={\fontsize{8}{9.5}\selectfont}}]
    \addplot[mainline, black]          {lgk(x,2074,1.00,16)};
    \addlegendentry{$\Delta n_{d}=0$}
    \addplot[thinline, black, dashed]  {lgk(x,2074-20,1.00,16)};
    \addlegendentry{$\Delta n_{d}=0.20$}
    % 【重叠修正，2026-09-30】$-20$ 原用 anchor=north east 挂在 (2066,1.22)，
    % 落在图例第二行 `$\Delta n_d=0.20$` 上（实测 `=0`×`−20`）。移到图例右侧空白处。
    \node[anchor=west, font={\fontsize{7}{8.5}\selectfont}]
      at (axis cs:2092,1.02) {$-20$};
  \nextgroupplot[ymin=1950, ymax=2090, ytick={1980,2020,2060},
                 % 必须显式覆盖 group 的 domain=1700:2200：本面板横轴是 0—0.6，
                 % 沿用 1700—2200 的 domain 会报 Dimension too large（实测）。
                 domain=0:0.6, xmin=0, xmax=0.6,
                 title={(c) 位移—电子密度关系},
                 title style={font={\fontsize{8}{9.5}\selectfont}}]
    \addplot[mainline, black, samples=61] {mtop(x)};
    \addlegendentry{顶位（模型）}
    \addplot[thinline, black, dashed, samples=61] {mbr(x)};
    \addlegendentry{桥位（模型）}
    \addplot[only marks, mark=*, mark size=1.5pt] coordinates {(0.09,2075)};
    \addplot[only marks, mark=square*, mark size=1.5pt] coordinates {(0,2049)};
    % 【重叠修正，2026-09-30】这里原有的两行方块注记（`• 组成位移 29 cm⁻¹` /
    % `∘ 覆盖度位移 9 cm⁻¹`）与图例 `桥位（模型）` 相撞：实测 6 对字形墨迹相交
    % （最大 `桥`×`移` 0.903 pt²）。该面板内两条斜线贯穿全幅，注记方块
    % （宽 29 mm、高两行）放在任何位置都与斜线或图例相交；而注记内容与图题
    % 完全重复（图题已给出两个位移数值），故删除，改在图题中指明标记点形状。
\end{groupplot}
\end{tikzpicture}
\caption{CO 探针谱带的模型计算：构型、电子密度与位移的关系}
\label{fig:ch07b-co-model}
\modelmark\ 三幅均为模型计算，不是实验谱。(a) 用洛伦兹型峰形示意顶位与
桥位两条谱带的分区，竖线标出 $2000\;\mathrm{cm^{-1}}$ 这一常用分界，
桥位峰按 $2.5$ 倍放大以便比较峰宽；(b) 只改电子密度、保持覆盖度与峰宽
不变，展示每个电子约 $100\;\mathrm{cm^{-1}}$ 的红移率下
$\Delta n_{d}=0.20$ 对应的 $-20\;\mathrm{cm^{-1}}$ 位移；(c) 横轴为模型设定
的位移标尺（单位相当于「每个电子的 $100\;\mathrm{cm^{-1}}$」），两条直线
是顶位与桥位的位移设定，两个标记点按本书实测量级放置：\textbf{实心圆点}为
组成位移 $29\;\mathrm{cm^{-1}}$、\textbf{实心方点}为覆盖度位移
$9\;\mathrm{cm^{-1}}$。读者应看的是：
(a) 中两峰间距远大于 (b) 中的电子效应位移，故构型归属不正确时，位移
判读会被几何效应完全掩盖；(c) 中两个标记点在横轴上的间距说明，要区分
电子效应与覆盖度效应，两者对应的位移必须在量级上先被分开。局限在于斜率
取自模型设定而非标定实验，绝对位置不用于定量。
\figsrc{峰位分界与两类位移的实测数值据原书 §7.3.2（书内 p.452—453 /
PDF p.471—472）；曲线形式为本节模型设定。}
\end{figure}
\end{crosscheck}

\begin{pitfall}
\textbf{把共吸附位移当纯电子效应。}共吸附分子占据部分吸附位、改变局部
覆盖度，本身就会带来位移；只有像本书那样把多个给电子能力不同的分子
排成序列、看\textbf{排序是否与给电子能力一致}，才能把电子效应从几何
与覆盖度效应中分出来。

\textbf{忽略共吸附分子的自身转化。}给电子分子可能与被测表面发生反应，
受主分子（如 $\ce{HCl}$）可能腐蚀或改性表面。谱带位移报告的是「体系
变化后的结果」，不必然是「该分子给电子」这一单一原因。

\textbf{把「位移量」当「转移电子数」。}图 \ref{fig:ch07b-co-model}(c)
中的直线是模型设定。真实体系中消光系数、偶极耦合与构型分布同时变化，
位移与电子转移量不是一一对应。
\end{pitfall}

\subsection{红外光谱在氧化物与分子筛催化剂表征中的应用}
\label{sec:ch07b-oxide}

\subsubsection{表面羟基的类型与酸强度}

\begin{probout}
要输出的是三件事：表面上存在\textbf{哪几种}结构羟基、每种羟基的\textbf{酸
强度}如何、以及各种羟基的\textbf{可接近程度}（能不能被探针分子接近）。
结构羟基的性质与表面酸性密切相关，而脱水反应、甲酸分解等许多催化反应
又与它直接相关\src{S1 §7.4.2，书内 p.465 / PDF p.484}，所以这三件事是
把红外谱与反应性能关联起来的必经环节。
\end{probout}

\begin{principle}
羟基的振动分析可近似为双中心的非谐振子模型。振动态表示为
\begin{equation}
  G(v)=\omega_{e}\!\left(v+\tfrac{1}{2}\right)
       -\omega_{e}x_{e}\!\left(v+\tfrac{1}{2}\right)^{2},
  \label{eq:ch07b-morse-levels}
\end{equation}
式中 $G(v)$ 为实验测得振动频率（以波数计），$\omega_{e}$ 为简正频率，
$x_{e}$ 为非谐振因子，$v$ 为振动能级数。
对应的 Morse 势能函数为
\begin{equation}
  U(r)=D\left[\,e^{-\beta(r-r_{e})}-1\,\right]^{2},
  \label{eq:ch07b-morse}
\end{equation}
式中 $r_{e}$ 与 $r$ 分别为振动原子的平均距离和瞬间距离，$D$ 为势阱深度，
$\beta$ 为与约化质量、简正频率有关的参数
\src{S1 §7.4.2，书内 p.470 / PDF p.489}。
\tocheck{未对照原页面图像：源 OCR 把 $\beta$ 与 $D$ 的表达式印成
「$\beta=2.48\times10^{7}(\cdots)$」「$D=\cdots$」等残形，本节只写出
可确证的标准 Morse 形式，$\beta$、$D$ 的具体表达式及其常数取值须核
原页 p.470。}
\end{principle}

\begin{principle}
用这套模型可以只靠一对同位素光谱算出解离能。由同一羟基的 OH 与 OD
伸缩振动基频（分别记 $v_{\mathrm{OH}}$、$v_{\mathrm{OD}}$，对应简正频率
$\omega_{\mathrm{H}}$、$\omega_{\mathrm{D}}$），可解出简正频率与非谐振
因子，进而得到质子解离能
\begin{equation}
  D_{0}=\frac{\omega_{e}^{2}}{4\,\omega_{e}x_{e}} .
  \label{eq:ch07b-D0}
\end{equation}
本书指出，羟基的质子解离能（即质子转移活化能）与 B 酸中心的酸性及固体
表面催化性能有密切联系；由吸附前后羟基与氘化羟基谱带的化学位移，可算出
有效质子转移活化能 $D_{\mathrm{eff}}$\src{S1 §7.4.2，书内 p.471 /
PDF p.490}。

\textbf{成立条件}：式 \equref{eq:ch07b-D0} 建立在「吸附分子的主要影响是
增大羟基的非谐振性、而简正频率基本不变」这一近似之上。本书明确指出，
只有在该近似下，才能从一个羟基基振动红外跃迁的化学位移求出质子转移
活化能的降低量；若吸附同时改变简正频率，则须用完整的式
\equref{eq:ch07b-morse-levels} 与 \equref{eq:ch07b-morse} 联立求解。
\tocheck{未对照原页面图像：本书由位移求 $D_{\mathrm{eff}}$ 的中间步骤
（源 OCR 残存「$D_{0}=\cdots/(\omega_{e}-2\omega_{e}x_{e})$」
「$X=(\omega_{e}-\omega_{e}^{*})/2\omega_{e}$」
「$\Delta\omega=\omega_{e}^{*}-\omega_{e}$」等形）无法从文本层可靠重建，
本节只保留式 \equref{eq:ch07b-D0} 与「由位移求活化能降低」这一结论，
中间式须核原页 p.471—472。}
\end{principle}

\begin{obs}
\textbf{二氧化硅表面。}在不同温度下脱水后测红外光谱，可发现吸附水对结构
羟基影响明显。特别的是，$800\,\degC$ 脱水后 $\ce{SiO2}$ 表面结构羟基
出现转动结构（$P$ 支、$R$ 支），表明羟基在 $\ce{SiO2}$ 表面上可以自由
转动，而在其他表面从未发现这一现象\src{S1 §7.4.2，书内 p.466—467 /
PDF p.485—486}。

\textbf{脱水温度与可接近程度。}$473\;\mathrm{K}$ 脱水后的硅溶胶表面，
$3740\;\mathrm{cm^{-1}}$ 处出现一尖谱带，宽谱带在 $3660\;\mathrm{cm^{-1}}$
与 $3550\;\mathrm{cm^{-1}}$；由同位素交换（OH $\to$ OD，
$2740\;\mathrm{cm^{-1}}$ 一带为 OD 伸缩带）结果看出，
$3550\;\mathrm{cm^{-1}}$ 峰与 $3660\;\mathrm{cm^{-1}}$ 峰是不可接近的
\src{S1 §7.4.2，书内 p.467 / PDF p.486}。
\tocheck{未对照原页面图像：源 OCR 作「$3740\;\mathrm{cm^{-1}}$ 峰、
$3550\;\mathrm{cm^{-1}}$、$3660\;\mathrm{cm^{-1}}$ 峰是不可接近的
（$3\cdots0\;\mathrm{cm^{-1}}$ OH 分 $27\cdots00$）」，可接近与不可接近
三种羟基的\textbf{一一对应关系}在文本层残缺，须核原页 p.467 与图 7-53。}

\textbf{氧化铝表面。}$\gamma\text{-}\ce{Al2O3}$ 经严格脱水后，红外光谱中
存在五个不同吸收带，对应五种不同的结构羟基。进一步研究发现各种结构
羟基的电子云密度顺序为 $\mathrm{A}>\mathrm{D}>\mathrm{B}>\mathrm{E}>\mathrm{C}$，
即 A 类部位最负、碱性最强，C 类最正、酸性最强，B 类部位近于中性
\src{S1 §7.4.2，书内 p.468 / PDF p.487}。

\textbf{羟基类型的一般图景。}把上述各体系合起来看，结构羟基可按其与骨架
的连接方式分成三类：\textbf{端式}（单个 $\mathrm{Si}$ 或 $\mathrm{Al}$
上的 $\equiv\!\mathrm{Si{-}OH}$）、\textbf{桥式}（$\equiv\!\mathrm{Si{-}(OH){-}Al}\equiv$
一类跨两个骨架原子的羟基，本书在分子筛处按变角振动单独列出）与
\textbf{巢式}（羟基位于骨架缺陷／空位处，配位环境不完整）。
\textbf{三者的酸强度顺序为端式 $<$ 巢式 $<$ 桥式}，因为桥式羟基所连的
$\mathrm{Al}$ 电负性低于 $\mathrm{Si}$、$\mathrm{Al{-}O}$ 键更长，
$\mathrm{O{-}H}$ 上的电子云密度被拉得更低，质子更易给出；巢式羟基因
配位不完整而处于两者之间。\textbf{请注意本节的证据边界}：源文本明确给出
并可用数据支持的是\textbf{端式与桥式}两类——$\ce{SiO2}$、无定形硅酸铝、
X 型、Y 型、丝光型分子筛都在 $3745\;\mathrm{cm^{-1}}$ 处有
$v_{\mathrm{OH}}$ 带（端式），而桥式羟基的变角振动与端式可相差
$180$—$250\;\mathrm{cm^{-1}}$\src{S1 §7.4.2，书内 p.472—473 /
PDF p.491—492}；「巢式」的命名与上述三者酸强度的\textbf{定量排序}，
源文本只以「可接近／不可接近」与电子云密度顺序间接涉及，未给出该排序
本身。\tocheck{未对照原页面图像：三羟基类型的命名（端式／桥式／巢式）
与「端式 $<$ 巢式 $<$ 桥式」这一酸强度排序，源文本 §7.4.2 未直接写出，
系按固体酸领域的通行分类补入；本书是否采用同一命名与排序须核
书内 p.466—469 / PDF p.485—488。}
\tocheck{未对照原页面图像：源 OCR 中表 7-7 的五种羟基波数只完整保留了
$3800$、$3780$、$3744$、$3733$、$3700$ 五个数，与其类型字母
（A—E）及最近邻配位数的\textbf{对应关系}已错行，此处不逐一列表，
须核原页 p.469。}
\end{obs}

\begin{infer}
由此得到「羟基类型—酸强度」的判读规则：\textbf{同一表面上的羟基并非
等价}，其酸强度由该羟基所处配位环境的电子云密度决定；电子云密度越低
（越缺电子），质子越易给出，酸强度越高。A 类最负故碱性最强、C 类最正
故酸性最强，中间的 B 类近中性，就是这条规则的直接体现。
判读时因此不能把「$3700$ 多波数区的一簇峰」当成一种羟基，也不能把
总羟基量当作酸量。

\textbf{氢键引起的位移。}当羟基与吸附质形成氢键时，羟基伸缩带向低波数
位移、同时强度增加、半宽增大。本书用氢受主化合物（$\ce{C6H6}$、
$\ce{CH3CN}$）吸附后的红外光谱来表征表面羟基的酸性强度：由键合作用引起
的羟基伸缩吸收带的\textbf{化学位移}可估计有关 OH 基的酸性，并可识别
表面 OH 的类型\src{S1 §7.4.2，书内 p.469 / PDF p.488}。
在含 $75\%$ $\ce{SiO2}$ 的 $\ce{SiO2{-}Al2O3}$ 上，吸附受主化合物后，
$3750\;\mathrm{cm^{-1}}$ 一带发展成两个峰，说明表面存在两种 OH：
一种构型为 $\equiv\!\mathrm{Si{-}OH}$，类似于纯 $\ce{SiO2}$ 表面；
另一种为 $\equiv\!\mathrm{Al{-}(OH){-}Si}\equiv$，其酸性较强，类似于
硅被 $\ce{Al^{3+}}$ 取代后的桥式羟基\src{S1 §7.4.2，书内 p.469 /
PDF p.488}。\textbf{位移的幅度即酸强度的量度}，这是把「位置」用于
强度判断的少数几处之一，成立前提是探针分子的碱性已知且同一批实验内
一致。
\end{infer}

\begin{crosscheck}
\textbf{用近红外组频带把「看不清的分辨开」。}由于强烈的本底，羟基变形
振动（$700$—$1100\;\mathrm{cm^{-1}}$）的信息过去一直无法获得。用近红外
方法考察分子筛的基频、倍频和组频后，可由羟基伸缩振动 $v_{\mathrm{OH}}$
与变角振动 $\delta_{\mathrm{OH}}$ 的组频 $(v+\delta)$、以及伸缩振动基频，
通过下式求出变角振动频率 $\delta_{\mathrm{OH}}$
\begin{equation}
  (v+\delta)=v_{\mathrm{OH}}+\delta_{\mathrm{OH}}-2X_{1},
  \label{eq:ch07b-combination}
\end{equation}
式中 $X_{1}$ 为两振动模式的相互作用系数，本书指出对分子筛而言它与
$\ce{SiO2}$ 的值相近\src{S1 §7.4.2，书内 p.472 / PDF p.491}。
\tocheck{未对照原页面图像：源 OCR 把式（7-11）印成
「嘿$=$嘿$+$嘿$+$嘿」一类的重复字形，$X_{1}$ 的系数（$2X_{1}$ 还是
$X_{1}$）及其数值量级在文本层无法判定；本节按「组频 $=$ 两个基频之和
减去相互作用项」的标准形式重建，须核原页 p.472。}

这一手段的价值在于\textbf{灵敏度}：$\ce{SiO2}$、无定形硅酸铝、X 型、
Y 型、丝光型分子筛都在 $3745\;\mathrm{cm^{-1}}$ 处有一个 $v_{\mathrm{OH}}$
吸收带，靠基频无法区分；但它们的 $\delta_{\mathrm{OH}}$ 明显不同，
桥式羟基与端式羟基的变角振动可相差 $180$—$250\;\mathrm{cm^{-1}}$。
更直接的例子是：X 型分子筛在桥式羟基基频区只有一个 $3650\;\mathrm{cm^{-1}}$
吸收带，而在组频带却分别在 $4650$ 与 $4690\;\mathrm{cm^{-1}}$ 出现两条，
表明该区存在两种羟基，但在基频范围无法分辨
\src{S1 §7.4.2，书内 p.472—473 / PDF p.491—492}。结论是：
\textbf{变角振动对羟基周围环境变化的灵敏度大于伸缩振动基频}，故基频
分不开时应当回查近红外组频带。
\tocheck{未对照原页面图像：源 OCR 中表 7-11 的端基／桥式羟基行大量损坏
（端基 $\mathrm{Si{-}OH}$ 的 $\delta$ 值残存 $795$、$820$、$805$、$825$
等数，桥式羟基行残存 $835$ 等数），各分子筛与各羟基类型的
$v+\delta$ 与 $\delta$ 一一对应关系须核原页 p.472—473。}

\textbf{用同位素交换独立验证。}$\ce{SiO2}$、$\ce{Al2O3}$、
$\ce{SiO2{-}Al2O3}$ 表面羟基与 $\ce{CD4}$ 的交换能力与酸性\textbf{没有}
直接关系——$\ce{SiO2}$ 上羟基中的质子最难交换
\src{S1 §7.4.2，书内 p.470 / PDF p.489}。这一条是对「酸性强则质子活泼」
这一直觉的反例：判断可接近性必须另做同位素实验，不能由酸强度外推。
\end{crosscheck}

\subsubsection{固体表面酸性（L 酸与 B 酸）的测定}

\begin{principle}
需要测定的是表面酸性位的\textbf{类型}（L 酸／B 酸）、\textbf{强度}与
\textbf{酸量}三者。碱滴定法、碱性气体吸附法、热差法等都不能区分 L 酸
部位与 B 酸部位，红外光谱法则可以有效区分
\src{S1 §7.4.3，书内 p.474 / PDF p.493}。
区分依据是探针分子在两类位上的\textbf{成键方式不同，因而环变形振动的
位置不同}：吡啶以氮的孤对电子配位到 L 酸部位，保持中性分子形式；在
B 酸部位则被质子化生成吡啶阳离子 $\mathrm{PyH^{+}}$。二者在
$1400$—$1700\;\mathrm{cm^{-1}}$ 的环变形振动区给出不同的谱带组合。
归属位置与抽空条件汇总于 \tabref{tab:ch07b-py}。
\end{principle}

\begin{table}[htbp]
\centering
\caption{吡啶吸附物种的环变形振动归属与抽空条件}
\label{tab:ch07b-py}
\small
\begin{tabularx}{\onecolwidth}{@{}>{\raggedright\arraybackslash}p{26mm}
 >{\raggedright\arraybackslash}p{34mm}
 >{\raggedright\arraybackslash}X@{}}
\toprule
吸附形式 & 归属谱带$^{\mathrm{a}}$ & 抽空行为与判读要点 \\
\midrule
物理吸附吡啶 &
 约 $1445$（$19b$）、约 $1490$（$19a$） &
 室温即可抽除；与氯仿溶液中的纯吡啶谱对应。测化学吸附物种时必须用
 $150\,\degC$ 抽空把它扣掉 \\
氢键吸附的吡啶 &
 约 $1450$（$19b$）、约 $1490$（$19a$） &
 $150\,\degC$ 可抽除。$19b$ 位置与 L 酸位配位吡啶几乎重合，
 不能单独用它下结论 \\
L 酸位配位吡啶 &
 \textbf{约 $1450$}（$19b$）、约 $1490$（$19a$） &
 室温抽空后保留，须 $150\,\degC$ 以上才逐步脱除；
 $1450$ 带是 L 酸位的判据带 \\
B 酸位质子化吡啶（$\mathrm{PyH^{+}}$） &
 \textbf{约 $1540$}（$19b$）、约 $1490$（$19a$） &
 强吸附，抽空至 $420\,\degC$ 仍可保留；$1540$ 带是 B 酸位的判据带。
 质子化产生的其他谱带位置受氢键影响变化大，不作定量依据 \\
两类位共有 &
 约 $1490$（$19a$） &
 \textbf{不能用它区分 L 酸与 B 酸}，只能用来说明「存在酸位」 \\
\bottomrule
\end{tabularx}

\vspace{1pt}
{\srcsize $^{\mathrm{a}}$ 单位 $\mathrm{cm^{-1}}$；$19b$、$19a$ 为环变形
振动的模式记号。数值为文献通行归属值，不同文献间可有约
$10\;\mathrm{cm^{-1}}$ 的差异，比对时须用同一套文献值。}
\tocheck{未对照原页面图像：本书表 7-12（书内 p.475 / PDF p.494）在源 OCR
中列间错位严重——$19b$ 与 $19a$ 两列的数被拆到不同行，「$1457$」
「$1625$」「$1615$」「$1540$」等数各自所属的吸附形式无法从文本层判定；
本表只列出源文本与正文别处可相互印证的带位，其余带位与全部
\textbf{最近邻配位数、峰半宽}数据一律不列，须核原页 p.475 后补齐。}
\end{table}

\begin{measure}
条件设置中三件事决定成败：
\textbf{(1) 探针碱性决定分辨能力。}本书指出，吡啶的碱性弱于 $\ce{NH3}$，
它能同弱酸部位反应\src{S1 §7.4.3，书内 p.474 / PDF p.493}；
反过来，$\ce{NH3}$ 碱性过强，会把强、弱酸位一并滴定，从而掩盖弱酸位的
差异。
\textbf{(2) 必须扣除物理吸附与气相。}为排除物理吸附吡啶的影响，应在
吸附后再于 $150\,\degC$ 下抽空，然后测定红外光谱
\src{S1 §7.4.3，书内 p.475 / PDF p.494}。物理吸附物种在室温即可抽除，
因而「室温吸附」与「$150\,\degC$ 抽空后」两张谱的差，就是化学吸附物种。
\textbf{(3) 选对用于定量的谱带。}吡啶分子面内环变形振动在若干位置出现，
但并非每条都适合定量（见下文误判）。
\end{measure}

\begin{obs}
\textbf{参比体系。}在氯仿溶液中：纯吡啶谱相当于物理吸附态；吡啶与典型
电子对受体 $\ce{BH3}$ 形成的 $\mathrm{C5H5N{+}{-}BH3}$ 相当于吸附在
L 酸部位；吡啶与 $\ce{HCl}$ 形成的 $\mathrm{C5H5N{+}HCl}$（即
$\mathrm{PyH^{+}}$ 与 $\mathrm{Cl^{-}}$）相当于吸附在 B 酸部位
\src{S1 §7.4.3，书内 p.474 / PDF p.493}。
\tocheck{未对照原页面图像：源 OCR 把该参比谱图 7-58 的谱带数值大量丢失
（仅残存 $1465$、$1585$、$1610$、$1640$ 等零散数与
「$1520\;\mathrm{cm^{-1}}$ 的宽峰由溶剂氯仿引起」一句），三个参比物种
各自的完整谱带列表须核原页 p.474 与图 7-58。}

\textbf{表面酸性测定的实例。}
$\ce{SiO2}$：吡啶在其上只是物理吸附，$150\,\degC$ 抽空后几乎全部脱附，
进一步表明纯 $\ce{SiO2}$ 上没有酸性中心。
$\ce{Al2O3}$：只能看到 L 酸中心（$1450\;\mathrm{cm^{-1}}$），看不到
B 酸中心。
$\ce{SiO2{-}Al2O3}$：在 $1600$—$1450\;\mathrm{cm^{-1}}$ 范围内除
$1450\;\mathrm{cm^{-1}}$ 的 L 酸带外，还出现 $1540\;\mathrm{cm^{-1}}$
的 B 酸带，说明该表面同时存在 L 酸部位与 B 酸部位；引入少量
$\ce{H2O}$ 后，吸附在 B 酸部位（$1540\;\mathrm{cm^{-1}}$）的吡啶浓度
增加，而 L 酸部位（$1450\;\mathrm{cm^{-1}}$）的浓度减少
\src{S1 §7.4.3，书内 p.475—476 / PDF p.494—495}。
由该体系还可读出四条更细的结论：表面至少存在三种 L 酸部位，其中一种
与混合相有关（加水可转变成 B 酸），另外两种一强一弱与 $\ce{Al2O3}$ 相
有关（加水不能转变成 B 酸）；随 $\ce{Al2O3}$ 含量增加，L 酸部位增加，
且主要是增加弱酸部位；混合相增加时强酸部位增加
\src{S1 §7.4.3，书内 p.476 / PDF p.495}。
\end{obs}

\begin{infer}
\textbf{L／B 归属的读法}：认定某样品有 B 酸位，依据是
$1540\;\mathrm{cm^{-1}}$ 带的出现，它对应质子化物种 $\mathrm{PyH^{+}}$
的环变形振动；认定有 L 酸位，依据是 $1450\;\mathrm{cm^{-1}}$ 带，它对应
以配位键结合的中性吡啶。必须同时看这两条带，因为它们不是同一条带的位移，
而是两个不同物种的谱带；只看一条带无法判断另一类酸位是否存在。
$1490\;\mathrm{cm^{-1}}$ 一带为两类位共有，因此\textbf{不能用它区分}
L 酸与 B 酸。

\textbf{「加水使 L 酸转成 B 酸」的含义}：$1540$ 带增强而 $1450$ 带减弱，
说明水把一部分 L 酸部位转化成了能给出质子的 B 酸部位，而不是简单地把
探针分子洗掉。这一现象把「酸位类型」与「预处理／含水状态」绑定在一起，
因此酸位类型的报告必须同时报告预处理条件。

\textbf{$\ce{NH3}$ 作为探针的归属}：$\ce{NH3}$ 吸附在 L 酸部位时以氮的
孤对电子配位，其红外光谱类似于金属离子与 $\ce{NH3}$ 的配位络合物；
在 B 酸部位则生成 $\mathrm{NH4^{+}}$。本书给出的区分主要依靠
$\delta_{\mathrm{N{-}H}}$ 变形振动：B 酸位与 L 酸位分别在
$1432\;\mathrm{cm^{-1}}$ 与 $1610\;\mathrm{cm^{-1}}$ 一带给出谱带
\src{S1 §7.4.3，书内 p.476—477 / PDF p.495—496}。注意该归属表引自
文献，不同文献的数值有约 $10\;\mathrm{cm^{-1}}$ 的差异，故比对时须用
同一套文献值。
\end{infer}

\begin{pitfall}
\textbf{误判一：把 $1490\;\mathrm{cm^{-1}}$ 带当 L 酸或 B 酸的判据。}
该带为两类位共有，用它判断酸位类型会把「有无酸位」错读成「有哪类酸位」。

\textbf{误判二：不扣物理吸附与气相就报酸量。}扣除的前提是有明确的抽空
温度阈值。本书的做法是吸附后在 $150\,\degC$ 抽空再测
\src{S1 §7.4.3，书内 p.475 / PDF p.494}。若省略这一步，室温可抽除的
物理吸附吡啶会计入 B 酸或 L 酸量；若抽空温度过高，部分弱酸位上化学
吸附的探针也会脱附，酸量被低估。此外，吡啶在强酸位上的质子化产物本身
就是被测物种，探针的自身反应与测量对象不可分离。

\textbf{误判三：用 $\nu(\mathrm{N{-}H})$ 带表征 B 酸量。}由于该带受
氢键等因素影响变化大、不易确定，本书明确指出一般不用它表征 B 酸部位
\src{S1 §7.4.3，书内 p.475 / PDF p.494}。

\textbf{误判四：跨样品直接比强度。}不同样品自支撑片的厚度与散射本底
不同，方程 \equref{eq:ch07b-beer} 中的 $l$ 不可比。跨样品比较须做厚度
归一或以内标带归一（见 \S\,\ref{sec:ch07b-oxide} 前文的说明）。

\textbf{误判五：用碱性过强的探针找弱酸位。}$\ce{NH3}$ 会把强、弱酸位
同时滴定，掩盖弱酸位的差异；要分辨弱酸位应改用碱性较弱的探针
（如吡啶），必要时再配合不同尺寸的探针区分孔内与外表（见下文）。
\end{pitfall}

\begin{crosscheck}
\textbf{用探针分子尺寸区分「孔内」与「外表面」。}酸中心的\textbf{空间
分布}是另一个必须回答的问题，靠的是不同截面积的探针分子。吡啶的临界
直径在 $6\;\mathrm{\text{\AA}}$ 左右，测得的是\textbf{总的}酸中心
（含介孔与微孔表面）；$2,6-$二甲基吡啶的临界直径更大，只能接近外表面
及孔口附近较大的中心。本书对系列样品（介孔硅铝与不同结晶度的 ZSM-5）
比较后给出：$2,6-$二甲基吡啶可接近的 B 酸中心数与外表面 B 酸所占
百分比顺序为 $\mathrm{Al{-}Meso}<\mathrm{UL{-}ZSM{-}5{-}2}
<\mathrm{UL{-}ZSM{-}5{-}4}<\mathrm{ZSM{-}5}$
\src{S1 §7.4.3，书内 p.481 / PDF p.500}。
\tocheck{未对照原页面图像：$2,6-$二甲基吡啶的临界直径数值在源 OCR 中
作「11.5 人$^{2}$」，量纲（面积还是直径）与数值均须核原页 p.481。}

\textbf{用羟基带与吡啶带的对应关系交叉验证。}分子筛上可同时看到两种
结构羟基，其中 $3612\;\mathrm{cm^{-1}}$ 一带的羟基对应于吡啶吸附在其上
产生的 B 酸 $1540\;\mathrm{cm^{-1}}$ 带；用大尺寸探针后该羟基带与
$1540$ 带同时消失，进一步证明此中心存在于微孔内
\src{S1 §7.4.3，书内 p.481 / PDF p.500}。这种「羟基带与探针带成对
消失」的验证，比单独看 $1540$ 带可靠。

\textbf{用改性实验验证归属。}用沸石骨架振动谱带作内标的计算机红外
差谱方法发现：P 改性 HZSM-5 时，P 的主要作用是「杀掉」大部分强的
B 酸中心（位于 $3611\;\mathrm{cm^{-1}}$ 位置的羟基）；Mg 改性除「杀掉」
大部分强 B 酸中心外，还形成一部分 L 酸中心
\src{S1 §7.4.3，书内 p.478 / PDF p.497}。把这一结果与 $\ce{NH3}$ 的
TPD（HZSM-5 表面存在强、弱两类酸性中心，P 和 Mg 改性后强酸部位大为
减少）以及 ESCA 的表面组成结果相对照，三条独立证据指向同一结论，这种
「红外差谱 $+$ TPD $+$ 表面组成」的组合是酸位表征的常规回查方式。

\textbf{补充的探针物性参照。}本书表 7-16 列出部分胺类探针的碱度
（$\mathrm{p}K_{a}$）与临界直径供选择探针时参考
\src{S1 §7.4.3，书内 p.480 / PDF p.499}。
\tocheck{未对照原页面图像：该表在源 OCR 中列错行严重——同一分子
$\mathrm{C6H5N}$（吡啶）出现两次且记为「$5.21$」与「$11.12$」两个
相差甚远的 $\mathrm{p}K_{a}$，$\mathrm{C6H5NH2}$ 记为「$0.46$」，
而氨与水溶液中氨的标准 $\mathrm{p}K_{a}$ 为 $9.25$ 而源文正文作
「5」。故本表数值\textbf{一律不引用}，须核原页 p.480 后方可使用。}
\end{crosscheck}

\subsubsection{氧化物表面氧物种与低碳烃活化}

\begin{principle}
甲烷是烃类中组成最简单、对称性最高、化学上最惰性的分子，在常温下与
催化剂表面相互作用时间非常短暂，很难直接观察到它在表面的活化过程。
另一方面，碱性氧化物表面存在一层稳定的碳酸盐，去除这一保护层很困难，
使表面氧物种的研究很少进行\src{S1 §7.4.4，书内 p.483 / PDF p.502}。
要突破这两点，实验上必须走「\textbf{先净化表面、再低温捕集}」的路线：
在 $1000\;\mathrm{K}$ 以上用纯 $\ce{O2}$ 长时间处理以去除表面碳酸盐
保护层，随后骤冷至 $173\;\mathrm{K}$ 进行氧物种和甲烷吸附态的红外
光谱研究，并与化学捕集技术、同位素交换技术结合
\src{S1 §7.4.4，书内 p.484 / PDF p.503}。

\textbf{判据的原理}：甲烷的 $v_{1}$ 对称伸缩振动在气相中是红外无活性的
（Raman 活性）。若分子以降低对称性的方式吸附在表面上，该振动会变为
红外活性并出现位移——\textbf{「对称性下降」由红外活性出现来判定，
「活化程度」由位移量来判定}。这是本节所有归属的共同基础。
\end{principle}

\begin{obs}
\textbf{氧物种。}经 $1000\;\mathrm{K}$ 纯 $\ce{O2}$ 长时间处理并抽高真空
得到洁净的 $\ce{CeO2}$ 表面，骤冷至 $210\;\mathrm{K}$ 吸附
$^{16}\ce{O2}$ 后的红外光谱中，$1128\;\mathrm{cm^{-1}}$ 谱带被归属为
$\ce{O2^{-}}$ 物种；再用 $^{18}\ce{O2}$ 作同位素实验，由同位素位移
进一步确认 $1065\;\mathrm{cm^{-1}}$ 谱带为 $\ce{O2^{2-}}$ 物种。
这一结果表明在新鲜的 $\ce{CeO2}$ 表面至少存在两种氧物种
\src{S1 §7.4.4，书内 p.485 / PDF p.504}。
\tocheck{未对照原页面图像：源 OCR 把 $1128\;\mathrm{cm^{-1}}$ 与
$1065\;\mathrm{cm^{-1}}$ 两条谱带各自归属的物种名印成「$^{16}\ce{O2^{-}}$」
与「$^{18}\ce{O2^{2-}}$」等混形，两个归属与两种氧物种（超氧
$\ce{O2^{-}}$ 与过氧 $\ce{O2^{2-}}$）的\textbf{一一对应关系}须核
原页 p.485 与图 7-71。}

\textbf{甲烷吸附态。}新鲜的 $\ce{CeO2}$ 表面上甲烷低温吸附后，在
$2875\;\mathrm{cm^{-1}}$ 出现吸收带；而气相 $\ce{CH4}$ 的
$v_{1}$ 在 $2917\;\mathrm{cm^{-1}}$ 处（Raman）并无红外活性。即：
甲烷吸附在 $\ce{CeO2}$ 表面导致 $\mathrm{C{-}H}$ 键振动由 Raman 活性
变为红外活性（对称性下降），并且与自由 $\ce{CH4}$ 分子相比存在
$42\;\mathrm{cm^{-1}}$ 的化学位移，表明甲烷分子被活化
\src{S1 §7.4.4，书内 p.485—486 / PDF p.504—505}。
不同氧化物上的位移量不同：$\ce{CeO2}$ 上约 $42\;\mathrm{cm^{-1}}$，
$\ce{MgO}$ 上 $17\;\mathrm{cm^{-1}}$，$\ce{Al2O3}$ 上
$27\;\mathrm{cm^{-1}}$\src{S1 §7.4.4，书内 p.486 / PDF p.505}。

\textbf{晶面对吸附的影响。}用 CO 探针研究 $\ce{CeO2}$ 单晶表面氧缺陷的
化学活性时发现：CO 在理想 $\ce{CeO2}(111)$ 晶面的红外振动峰位于
$2154\;\mathrm{cm^{-1}}$，而在还原表面的峰位于 $2163\;\mathrm{cm^{-1}}$；
CO 在缺陷位的红外峰显示出很小的波数变化，远远小于原来的预期，且
\textbf{不同于 $\ce{TiO2}$ 等材料上缺陷位吸附红移的现象}。借助理论计算
归属后认为，$2154\;\mathrm{cm^{-1}}$ 的峰来源于紧邻缺陷的六配位
$\ce{Ce^{3+}}$ 上吸附的 CO，而非离氧缺陷更远处的七配位 $\ce{Ce^{3+}}$
\src{S1 §7.4.5，书内 p.486—487 / PDF p.505—506}。
\end{obs}

\begin{infer}
\textbf{位移—活化程度的读法}：$42\;\mathrm{cm^{-1}}$（$\ce{CeO2}$）
$>27\;\mathrm{cm^{-1}}$（$\ce{Al2O3}$）$>17\;\mathrm{cm^{-1}}$
（$\ce{MgO}$）这一顺序，说明同一探针分子（甲烷）在不同氧化物表面上被
活化的程度不同；位移越大表示 $\mathrm{C{-}H}$ 键被削弱得越多。因此
\textbf{「位移量」在这里被用作活化程度的相对标尺}，比较必须在同一探针、
同一吸附温度下进行。

\textbf{配位数而非距离决定吸附位置}：$\ce{CeO2}$ 上 CO 在缺陷位的位移
\textbf{很小}，与 $\ce{TiO2}$ 等材料上缺陷位吸附\textbf{红移}的常见
现象相反。理论计算给出的解释是：该峰来自紧邻缺陷的六配位
$\ce{Ce^{3+}}$，而非更远处的七配位 $\ce{Ce^{3+}}$。由此得到的判读
规则是——\textbf{氧化物的配位结构（配位数）在表面化学相互作用中起
决定性作用}，不能仅凭「离缺陷的距离」推断吸附位与位移方向
\src{S1 §7.4.5，书内 p.486—487 / PDF p.505—506}。这一结论与本书对
甲氧基的研究相互印证：甲氧基 $\mathrm{C{-}O}$ 键伸缩振动峰的红移与
甲氧基氧原子配位的 Ce 阳离子数\textbf{不存在直接关系}
\src{S1 §7.4.5，书内 p.487 / PDF p.506}——即不能用「配位数越多红移越大」
这类单因素规则去套。
\end{infer}

\begin{pitfall}
\textbf{误判一：把位移单因素归因于电子效应或配位数。}本书用两处实测
反例说明这一点：$\ce{CeO2}$ 缺陷位 CO 的位移小于预期且方向与
$\ce{TiO2}$ 相反；甲氧基 $\mathrm{C{-}O}$ 红移与配位 Ce 阳离子数无直接
关系。凡遇到「位移不符合预期」的情形，正确做法是补做同位素实验与理论
计算，而不是修改归属规则去迁就数据。

\textbf{误判二：不区分「对称性下降」与「化学位移」两个信息来源。}
$v_{1}$ 由 Raman 活性变为红外活性只证明吸附降低了对称性；键被削弱的
程度必须另由位移量判断。把二者混为一谈会把「分子被表面固定住」误读成
「分子被活化」。

\textbf{误判三：忽略表面碳酸盐保护层。}碱性氧化物表面天然覆盖一层稳定
碳酸盐，不去除就直接做吸附实验，测到的主要是碳酸盐的变化而不是氧物种
与烃的相互作用。去除条件（$1000\;\mathrm{K}$ 以上纯 $\ce{O2}$ 长时间
处理）本身就是这套实验的关键条件，不能省略。

\textbf{误判四：把低温吸附态当成反应条件下的真实中间物。}$173\;\mathrm{K}$
骤冷下捕集到的吸附态是\textbf{稳定化}的观察对象；从它推断反应条件下的
活化路径，需要另有证据（如升温过程中的谱带演化、同位素标记实验）。
\end{pitfall}

\begin{crosscheck}
\textbf{同位素交换是归属的第一道回查。}氧物种的归属靠
$^{16}\ce{O2}/^{18}\ce{O2}$ 的同位素位移确认；甲烷吸附态的归属靠
$\ce{CD4}$ 同位素实验与同 CO 共吸附的双分子探针方法证实
\src{S1 §7.4.4，书内 p.486 / PDF p.505}。凡涉及含氧或含氢物种的归属，
同位素位移应当与位置归属同时报告。

\textbf{双分子探针是电子转移方向的回查。}用 CO 共吸附观察
$v_{\mathrm{CO}}$ 的位移方向（见 \S\,7.3.3 的读图实例），可以判断吸附
过程中电子转移的方向与程度，这一手段与氧物种的红外归属互为补充
\src{S1 §7.3.3，书内 p.457 / PDF p.476}。

\textbf{分子筛体系的回查逻辑可以照搬。}原位红外与拉曼光谱结合、X 射线
衍射与红外光谱结合，是把「纯相的谱」与「实际催化剂的谱」对应起来的
通行做法：本书用柠檬酸法在较低焙烧温度下获得较大比表面的
V—Mg—O 三种纯相（偏钒酸镁 $\ce{MgV2O6}$、正钒酸镁 $\ce{Mg3V2O8}$、
$\alpha\text{-}\ce{Mg2V2O7}$），由 X 射线衍射、红外光谱和拉曼光谱三者
的对应关系分辨纯相，再与实际催化剂的结果相对照
\src{S1 §7.4.1，书内 p.462—463 / PDF p.481—482}。同理，分子筛骨架振动
可用于测 Si/Al 值并对结构变化敏感——随 Si/Al 值减小，力常数减小、
波数降低；约 $1230\;\mathrm{cm^{-1}}$ 与约 $550\;\mathrm{cm^{-1}}$ 的
谱带对 ZSM-5 的结构变化敏感，约 $1230\;\mathrm{cm^{-1}}$ 谱带的化学
位移同 ZSM-5 的 Si/Al 比呈直线关系，这给出一条用红外测定分子筛 Si/Al
值的途径\src{S1 §7.4.1，书内 p.464—465 / PDF p.483—484}。
\textbf{回查规则}：凡要用一条谱带的位移去标定组成或结构参数，必须先有
X 射线衍射等独立手段确认「结构确实按该参数单调变化」。
\tocheck{未对照原页面图像：源 OCR 中表 7-6（焙烧温度对 ZSM-5 骨架振动
的影响）与图 7-49（$1230\;\mathrm{cm^{-1}}$ 谱带位置与 Si/Al 比的关系）
的数据格全部为空或残形，斜率与拟合式无法读出；本段只保留「成直线关系」
这一定性结论，定量关系须核原页 p.464—465。}
\end{crosscheck}

\begin{srcnote}
\S 7.3 据原书 §7.3（书内 p.448—458 / PDF p.467—477）：探针分子的判读
链条、双探针法测表面组成、几何效应与电子效应、吸附分子相互作用；
\S 7.4 据原书 §7.4.2—7.4.4（书内 p.465—488 / PDF p.484—507）：
表面羟基、固体表面酸性、氧化物表面氧物种与低碳烃活化，另从 §7.4.1
（书内 p.462—465 / PDF p.481—484）取分子筛骨架振动一条用于回查。
表 \ref{tab:ch07b-probes} 的谱带位置取自原书 §7.4.3 的归属
（书内 p.475、p.477 / PDF p.494、p.496）与 §7.4.2 的羟基与碳酸盐数据；
吡啶、$\ce{NH3}$ 的位置为文献归属值，不同文献间有约
$10\;\mathrm{cm^{-1}}$ 的差异。图 \ref{fig:ch07b-co-model} 为本节模型
计算，不是原书图，其峰位分界与两类位移的实测数值取自 §7.3.2。
原书 §7.2「吸附分子的特征和它的红外光谱诊释」属分片 a 范围，本节只在
判读链条处引用其结论。凡源 OCR 失真、未能对照原页面图像重建的公式、
数值与一一对应关系，已就地用 \tocheck{} 标出。
\end{srcnote}
